\section{The author's unpublished notes, April--June 2026}\label{sec:notes}

At the author's request, this version also reads the author's unpublished notes on the equation. They were written from April to June 2026, with help from language models, and are not part of the workbench's public documents. There are two collections:
\begin{itemize}
\item one of twenty \TeX{} notes and one 52-page paper;
\item the Erdős--Straus part of a second collection, about fifty-five further notes, among them the source note of the ``terminal core''.
\end{itemize}
They are cited here by title and theorem number.

The notes were read in six independent passes by separate Claude instances. Each pass tabulated every numbered statement and checked it by its own programs. Every result proved in this section was then re-derived and re-checked with this reader's own programs (Appendix~\ref{app:verification}). The statements that rest only on a reading pass, on the referee of this version, or on a cited source say so. The folder of drafts that the author labelled as overstating results was not opened. The full record is note \note{50\_}.

\paragraph{What the notes contain.}
Their Erdős--Straus arithmetic consists of the shell calculus of Section~\ref{sec:core} and the classical identity families of Rosati, Yamamoto, Elsholtz--Tao, Salez and Lopez.

Their structural layers are correct where they were checked:
\begin{itemize}
\item the Niemeier lattice $A_5^4D_4$, which is also archive Theorem~25.8, and its umbral group $GL_2(3)$, the automorphism group of the glue code (checked by a reading pass);
\item the Hauptmodul of $\Gamma_0(6)$ and the Monster class $6B$;
\item Ogg's primes;
\item Lorentzian and Apollonian geometry;
\item Cayley--Dickson algebras;
\item the split-zero semiring;
\item mock theta functions.
\end{itemize}
None of these structures takes a prime, a shell or a certificate as input, so none can constrain which primes are solvable. The notes mostly say so themselves.

\paragraph{What they add.}
This section records:
\begin{itemize}
\item what the notes add to Sections~\ref{sec:core} and~\ref{sec:families};
\item one correction of scope, namely that the ``terminal core'' is a support statement;
\item a short list of errors and connections.
\end{itemize}

\subsection{The terminal core is a support statement}\label{ssec:termcore}

Put $M_{23}=840\cdot11\cdot19\cdot23=4{,}037{,}880$. A \emph{base class} is a class $n\equiv r\pmod{840}$ with $r\in S_{840}=\{1,121,169,289,361,529\}$, together with units modulo $11$, $19$ and $23$. There are $23{,}760$ base classes.

The source note, \emph{Finite central-cover reductions for residual Erdős--Straus shells}, calls $(R,m,u)$ an \emph{elementary $j=1$ certificate} when:
\begin{itemize}
\item $R\equiv3\pmod4$ and $u\mid m^2$;
\item $n\equiv-R\pmod{4m}$, $n\equiv-4u\pmod R$, and $\gcd(n,R)=1$.
\end{itemize}
Then $m\mid a=(n+R)/4$, $u\mid a^2$ and $u\equiv-a\pmod R$, so the shell is occupied through the middle channel $M_a$. The primes that may divide $m$ are those of $\{2,3,5,7,11,19,23\}$ not dividing $R$ whose divisibility of $a$ the class forces.

The note makes two computations:
\begin{enumerate}[label=(\roman*)]
\item The $32$ shells $R\mid M_{23}$ with $R\equiv3\pmod4$ cover all but $4951$ base classes. This is a covering statement.
\item Over all $R$ built from $3,5,7,11,19,23$, it tests the certificate congruence modulo $\gcd(M_{23},R)$, so that a class counts as hit when some lift of it modulo $\operatorname{lcm}(M_{23},R)$ is certified. Exactly $2970$ classes are then never hit, namely $S_{840}\times Q_{11}\times Q_{19}\times Q_{23}$, where $Q_q$ is the group of non-zero squares modulo $q$.
\end{enumerate}
The note calls this set the ``support-zero core''. It uses the set to bound from below the residue of any obstruction built from such certificates.

\begin{lemma}\label{lem:j1square}
Let $(R,m,u)$ be an elementary $j=1$ certificate with $\gcd(m,R)=1$, and let $n$ satisfy its congruences. Put $\varepsilon_\ell(n)=(n/\ell)$ for odd primes $\ell$. Put $\varepsilon_2(n)=1$ if $n\equiv1\pmod8$ and $\varepsilon_2(n)=-1$ if $n\equiv5\pmod8$. Then the Jacobi symbol satisfies
\[
\Big(\frac nR\Big)=-\prod_{\ell\mid u}\varepsilon_\ell(n)^{v_\ell(u)}.
\]
In particular no $n$ in the class is a unit square modulo $\operatorname{lcm}(4m,R)$.
\end{lemma}
\begin{proof}
From $n\equiv-4u\pmod R$ we get $(n/R)=(-1/R)(4/R)(u/R)=-(u/R)$.

For an odd prime $\ell\mid m$, reciprocity with $R\equiv3\pmod4$ gives $(\ell/R)=(-R/\ell)$. This equals $(n/\ell)$ because $n\equiv-R\pmod\ell$.

If $2\mid m$, then $n\equiv-R\pmod8$. So $(2/R)=1$ exactly when $R\equiv7\pmod 8$, that is, when $n\equiv1\pmod8$.

If $n$ were a square modulo $\operatorname{lcm}(4m,R)$, every $\varepsilon_\ell(n)$ with $\ell\mid m$ would be $1$. Then $(n/R)=-1$, which is impossible for a square.
\end{proof}

\begin{theorem}[the terminal core]\label{thm:termcore}
\begin{enumerate}[label=(\alph*)]
\item A base class has no elementary $j=1$ certificate with $R$ built from $3,5,7,11,19,23$ on any lift if and only if it is a square modulo $11$, $19$ and $23$. So the support-zero core is exactly $S_{840}\times Q_{11}\times Q_{19}\times Q_{23}$, which has $2970$ classes.
\item The base classes that no single identity among Salez's seven modular equations \cite{Salez} with modulus dividing $M_{23}$ covers number $3520$. They are the $2970$ square classes and $550$ classes that are non-squares modulo $11$, $19$ or $23$.
\item So the classes left open at level $M_{23}$ strictly contain the core. For example the prime $3361\equiv1\pmod{840}$, $\equiv6\pmod{11}$, lies outside the core, and none of the $64{,}149$ families with modulus dividing $M_{23}$ covers its class. The primes $18{,}481$ and $2{,}840{,}041$ behave in the same way.
\end{enumerate}
\end{theorem}
\status{(a) The ``if'' direction is proved here. The ``only if'' direction is a finite computation: the source note's script and an independent program of this reader both give $4951$ and then $2970$. (b) A finite computation, calibrated at Salez's modulus $G_7=840\cdot11\cdot13\cdot17\cdot19\cdot23$, where it leaves $147{,}348$ classes, the number Salez reports. Salez states that his seven equations form a complete set of modular equations. The primes in (c) are of course solvable, for instance $4/3361=1/841+1/110139970+1/950330$.}
\begin{proof}
(a), ``if''. Let $n$ lie over a square class.
\begin{itemize}
\item $S_{840}$ is the group of unit squares modulo $840$, so $n$ is a quadratic residue modulo every prime $\ell\in\{3,5,7,11,19,23\}$.
\item Also $n\equiv1\pmod8$.
\item So $(n/R)=1$ for every $R$ built from these primes.
\item The primes that may divide $m$ lie in $\{2,3,5,7,11,19,23\}$, so Lemma~\ref{lem:j1square} forces $(n/R)=-1$ for any certificate.
\end{itemize}
Hence no lift is certified. The remaining statements are the computations in the status line.
\end{proof}

\begin{remark}\label{rem:termcore}
Two later papers treat the core as the unresolved set of hard classes: the 52-page \emph{Terminal residual coordinates for the Erdős--Straus problem} (26 May 2026), and its June replacement \emph{Finite support algebra, cubic middle geometry, and orientation descent}. At level $M_{23}$ that set is the $3520$-class set of Theorem~\ref{thm:termcore}(b).

In the vocabulary of the author's split-zero semiring (§\ref{ssec:archive}), used here only as an analogy:
\begin{itemize}
\item the $2970$ core classes are \emph{unsupported};
\item the $550$ further classes are \emph{supported but not covered}. Some lift of each is certified, but no identity of Salez's seven families holds on the whole class.
\end{itemize}
Salez's statement that his seven equations are complete concerns identity families of his linear type; the archive's audit of Salez (§13.2) records this scope.

By Dirichlet's theorem the core carries $2970/\varphi(M_{23})=1/256$ of all primes, and the classes left open at level $M_{23}$ carry $3520/760{,}320=1/216$.
\end{remark}

\subsection{Quadratic-residue conditions on a counterexample}\label{ssec:qrconditions}

Throughout this subsection, $p\equiv1\pmod{24}$ is prime. A certificate is a shell $R\equiv3\pmod4$ with $\gcd(R,pa)=1$ and a divisor $u\mid a^2$ with $R\mid4u+1$ (channel $E$) or $R\mid u+a$ (channel $M$). Theorem~\ref{thm:shell} then gives the solution.

\begin{theorem}[eight shifts]\label{thm:eightshifts}
Let $q$ be an odd prime dividing
\[
F(p)=(p+1)(p+2)(p+3)(p+4)(p+7)(p+8)(2p+1)(3p+1)
\]
with $(p/q)=-1$. Then the shell and divisor of Table~\ref{tab:eightshifts} give a certificate. Equivalently, a counterexample $p\equiv1\pmod{24}$ is a quadratic residue modulo every odd prime factor of $F(p)$.
\end{theorem}

\begin{center}\small
\begin{longtable}{@{}lllll@{}}
\caption{The certificates of Theorem~\ref{thm:eightshifts}, with $a=(p+R)/4$.}\label{tab:eightshifts}\\
\toprule
shift & $q$ with $(p/q)=-1$ & shell $R$ & $u$ & channel\\
\midrule
$p+1$ & $q\equiv3\pmod4$ & $q$ & $a$ & $E$\\
$p+4$ & $q\equiv3\pmod4$ & $q$ & $1$ & $M$\\
$p+2$ & $q\equiv5,7\pmod8$ & $q$ if $q\equiv7$, $3q$ if $q\equiv5\pmod 8$ & $a/2$ & $E$\\
$p+8$ & $q\equiv5,7\pmod8$ & the same & $2$ & $M$\\
$2p+1$ & $q\equiv5,7\pmod8$ & the same & $2a$ & $E$\\
$p+3$ & $q\equiv2\pmod3$ & $3$ & $q$ & $E$\\
$p+7$ & $q\equiv5$, $6$, $3\pmod7$ & $7$ & $q$, $aq$, $2aq$ & $E$, $M$, $M$\\
$3p+1$ & $q\equiv2\pmod3$ & $(4q+1)/3$ & $q$ & $E$\\
\bottomrule
\end{longtable}
\end{center}

\status{Proved here.
\begin{itemize}
\item \emph{Recorded rows.} $p+1$ and $p+4$ are archive Theorem~24.50. $p+3$ and $p+7$ are Propositions~\ref{prop:r3} and~\ref{prop:r7}, read through reciprocity. $3p+1$ is archive Theorem~24.171 with $k=3$.
\item \emph{Salez's charts.} Each of the rows $p+1$, $p+2$, $2p+1$, $p+4$, $p+8$ is the union, over the free divisor $R$, of one of Salez's charts \cite{Salez} with small fixed constants: (15b) with $(B,C)=(1,1),(1,2),(2,1)$ and (14c) with $(B,D)=(1,1),(1,2)$.
\item \emph{Not in the archive.} The quadratic-residue conditions from $p+2$ and $2p+1$, and the stated condition from $p+8$, are not stated in the archive, in this reader's earlier versions, or in \note{43\_}/\note{43a\_}. No literature search was made.
\item \emph{Checks.} Every certificate for every prime $p\equiv1\pmod{24}$ below $10^7$ was built and verified: $780{,}700$ of them, none bad. $247$ primes survive all eight shifts.
\end{itemize}}
\begin{proof}
\emph{The characters.} For $q\mid p+k$ with $q\nmid k$ we have $(p/q)=(-k/q)$; similarly $(p/q)=(-2/q)$ for $q\mid 2p+1$ and $(p/q)=(-3/q)$ for $q\mid3p+1$. These characters are $-1$ exactly on the classes of $q$ listed in the table; for $p+7$ this uses $(-7/q)=(q/7)$.

\emph{Each row is a certificate.} For each row the stated divisibility is an identity:
\begin{itemize}
\item $4a+1=p+R+1$;
\item $4(a+1)=p+R+4$;
\item $2(2a+1)=p+R+2$;
\item $4(a+2)=p+R+8$;
\item $8a+1=2(p+R)+1$;
\item $4q+1\equiv q+1\pmod3$;
\item for $p+7$: $4q+1\equiv0$, $a(q+1)$ and $a(2q+1)$ modulo $7$.
\end{itemize}

\emph{The rows $p+2$, $p+8$ and $2p+1$.} Since $p\equiv1\pmod3$, $3$ divides each of these shifts, so $R$ divides the shift. Here $q\neq3$ because $(p/3)=1$. Also $R\equiv7\pmod8$, so $8\mid p+R$ and $a$ is even.

\emph{The row $p+7$.} $8\mid p+7$, so $a$ is even, and $q\mid a$.

\emph{The row $3p+1$.} Put $N_3=(3p+1)/4$, which is odd and $\equiv1\pmod3$. Then $h=N_3/q\equiv2\pmod3$. With $a=q(h+1)/3$ one checks $4a-p=(4q+1)/3=R$.

\emph{Coprimality.} A common prime of $R$ and $a$ divides $4a-R=p$, and $p\nmid R$ in every row.
\end{proof}

At the strict record $p_*=8{,}803{,}369$, whose least occupied shell is $107$, two shifts give four certificates:
\begin{itemize}
\item $p_*+2=3\cdot223\cdot13159$, where $223\equiv7$ and $13159\equiv7\pmod8$, giving the shells $223$ and $13{,}159$;
\item $2p_*+1=3\cdot677\cdot8669$, where $677\equiv5$ and $8669\equiv5\pmod8$, giving the shells $2031$ and $26{,}007$.
\end{itemize}
None of the other six shifts gives a certificate at $p_*$. The least primes certified by $p+2$ alone, by $2p+1$ alone and by $p+8$ alone are $2521$, $9601$ and $33{,}961$.

\begin{theorem}[$p+5$]\label{thm:pplus5}
Let $N=(p+5)/2$, and let $p\equiv9\pmod{40}$, or $p\equiv121\pmod{840}$, or $p\equiv1,361\pmod{840}$ with $p\equiv4\pmod9$. Suppose some prime $q\mid p+5$ has $(p/q)=-1$. Then $N$ has a divisor $d\equiv-p\pmod{20}$ or a divisor $h\equiv-1\pmod{20}$, and either gives a certificate by Proposition~\ref{prop:locks} with $\mu=5$.

So a counterexample in these classes, which contain the hard classes $121$, $169$, $289$ and $529$, is a quadratic residue modulo every odd prime factor of $p+5$.
\end{theorem}
\status{Proved here. The case $p\equiv9\pmod{40}$ is in the note \emph{Full divisor locks, QR subtori, and survivor idempotents} in substance (its Theorem~3.7(ii), Corollary~3.8(ii) and Theorem~4.2). The extension to the classes $121$, $1$ and $361$ was found by a reading pass. The theorem was checked on all $41{,}419$ primes $p\equiv9\pmod{40}$ below $10^7$ and on $5688$ primes of the other classes. At $p=12601\equiv1\pmod{840}$, with $p\equiv1\pmod9$, it fails: $p+5=2\cdot3\cdot11\cdot191$ has two non-residue factors, yet $N$ has no divisor in either class.}
\begin{proof}
For $q\mid N$ we have $(p/q)=(-5/q)$. This equals $1$ exactly when $q\equiv1,3,7,9\pmod{20}$, a subgroup $H$ of index $2$.

\emph{The case $p\equiv9\pmod{40}$.} Here $N\equiv7\pmod{20}$, and the targets are $11$ and $19$. Let $q\notin H$ divide $N$:
\begin{itemize}
\item if $q\equiv11$ or $19$, then $q$ itself is the divisor;
\item if $q\equiv13$, then $N/q\equiv19$;
\item if $q\equiv17$, then $N/q\equiv11$.
\end{itemize}

\emph{The case $p\equiv1\pmod{40}$.} Here $N\equiv3\pmod{20}$, and a divisor $\equiv17$ or $19$ suffices, because its cofactor is then $\equiv19$. Let $q\notin H$ divide $N$:
\begin{itemize}
\item $3\mid N$ always, so $q\equiv13$ gives $3q\equiv19$;
\item $7\mid N$ when $p\equiv121\pmod{840}$, so $q\equiv11$ gives $7q\equiv17$;
\item $9\mid N$ when $p\equiv4\pmod9$, so $q\equiv11$ gives $9q\equiv19$.
\end{itemize}

\emph{Conversely}, the targets lie outside $H$, so a divisor in a target class has a prime factor outside $H$.
\end{proof}

\begin{theorem}[the visible-box principle]\label{thm:visiblebox}
Let $R\equiv3\pmod4$ be prime and $c$ a class of hard primes modulo $10080=2^5\cdot3^2\cdot5\cdot7$. Let
\[
\mathrm{Box}(c)=\Big\{\prod_{f\in\{2,3,5,7\}}f^{i_f}\bmod R:\ 0\le i_f\le2e_f\Big\},
\]
where $e_f$ is the least value of $v_f((p+R)/4)$ that $c$ forces. If $\mathrm{Box}(c)$ contains every non-zero square modulo $R$, then for every prime $p\in c$ the shell $R$ is occupied exactly when some odd prime $q\mid p+R$ has $(p/q)=-1$.
\end{theorem}
\status{Proved here; found by a reading pass. The archive's census refinements (Theorems~24.228 and~24.233) already give certificates in the shells $R=31$, $47$ and $71$ on parts of these classes; the equivalence is not stated there. The search up to $R=700$ is complete, because $|\mathrm{Box}(c)|\le7\cdot5\cdot3\cdot3=315$. Among the $72$ hard classes modulo $10080$ it finds $72$ good classes for $R=3$ and $R=7$, $48$ for $R=11$, $12$ for $R=19$, $48$ for $R=23$, and $30$, $20$, $4$, $12$, $1$ for $R=31$, $47$, $59$, $71$, $311$. The last five shells are new. All $90{,}897$ tests with $p<10^7$ agree, with $50{,}545$ certificates verified.}
\begin{proof}
Every element of $\mathrm{Box}(c)$ is the residue of a divisor of $a^2$.

If $q\mid a$ is a prime with $(q/R)=-1$, then $(p/q)=(-R/q)=(q/R)=-1$. Also $q\cdot\mathrm{Box}(c)$ is the non-square coset, so the divisors reach every unit, including the target $-4^{-1}$ of channel $E$.

The condition ``$(p/R)=-1$'' needs no separate treatment. Let $m$ be the odd part of $p+R$. Since $p\equiv1\pmod8$, reciprocity gives $(p/m)=(m/p)=(R/p)(2/p)^{-v_2(p+R)}=(p/R)$. So if $(p/R)=-1$, some prime factor $q$ of $m$ has $(p/q)=-1$.

Conversely, suppose that no odd prime $q\mid p+R$ has $(p/q)=-1$. For an odd prime $\ell\mid a$, reciprocity gives $(\ell/R)=(-R/\ell)=(p/\ell)=1$; the prime $2$ divides $a$ only when $R\equiv7\pmod8$, and then $(2/R)=1$. So every prime factor of $a$ is a square modulo $R$, and Proposition~\ref{prop:jacobi} shows that the shell is empty.
\end{proof}

\begin{proposition}[$\mu=17$]\label{prop:mu17}
Let $p\equiv361$ or $529\pmod{840}$, $p\equiv1\pmod9$ and $p\bmod17\in\{8,15,16\}$. Then $(p+17)/2$ has a divisor $d\equiv-p$ or $h\equiv-1\pmod{68}$, giving a certificate by Proposition~\ref{prop:locks}, if and only if some prime $q\mid p+17$ has $(p/q)=-1$.
\end{proposition}
\status{Proved here; found by a referee pass. Checked on all $433$ such primes below $10^7$.}
\begin{proof}
Put $N=(p+17)/2$. Then $63\mid N$. The character $(p/q)=(-1/q)(q/17)$ lives modulo $68$, and both targets $-p$ and $-1$ lie outside its kernel. A finite check of the six classes shows that for every non-kernel class $r$ the sets $r\cdot W$ and $N(rW)^{-1}$, with $W=\{1,3,7,9,21,63\}$, meet the targets modulo $68$. The converse holds as in Theorem~\ref{thm:pplus5}.
\end{proof}

Each of these conditions removes further residue classes modulo about half of all primes. The primes $p\equiv1\pmod{24}$ that satisfy them all are counted in Table~\ref{tab:survivors}.

\begin{center}\small
\begin{longtable}{@{}lrrr@{}}
\caption{Primes $p\equiv1\pmod{24}$ satisfying all the conditions of this subsection. The column $10^8$ is from a referee pass.}\label{tab:survivors}\\
\toprule
conditions & below $10^6$ & below $10^7$ & below $10^8$\\
\midrule
hard classes & $2370$ & $20{,}513$ & $179{,}468$\\
and the eight shifts & $54$ & $247$ & $1431$\\
and $p+5$ on its classes & $29$ & $146$ & $816$\\
and $\mu=17$ & $29$ & $145$ & $806$\\
and the visible-box shells & $7$ & $31$ & $165$\\
\bottomrule
\end{longtable}
\end{center}

The first primes that satisfy them all are $43{,}201$, $65{,}521$ and $196{,}561$.

The referee of this version found eight more conditions of the same kind, on parts of the hard classes. They come from $p+6$, $p+12$, $p+16$, $p+20$, $p+24$, $p+36$, $5p+1$ and $6p+1$, through the three certificate shapes $u=s$ with $R\mid p+4s$, $u=a/t$ with $R\mid p+t$, and $u=ta$ with $R\mid tp+1$. It verified $46{,}441$ certificates below $10^7$, and its script was re-run here. With these conditions, $23$ primes survive below $10^7$. The prime $43{,}201$ is removed by $p+24$: $4/43201=1/10914+1/1036824+1/1885982856$.

A sieve bound of the form $x/(\log x)^\kappa$ follows. It is far weaker than Vaughan's bound on the exceptional set \cite{Vaughan}, so these conditions measure what the explicit certificates sieve and do not approach a proof.

\begin{remark}[the Collatz reading]\label{rem:collatz}
For $p\equiv1\pmod8$ the number $(3p+1)/4$ is the odd Collatz successor of $p$. So a counterexample $p\equiv1\pmod{24}$ has an odd Collatz successor built only from primes $\equiv1\pmod3$.

The odd Collatz predecessors of an odd $n$ with $3\nmid n$ satisfy $\Pi_{k+1}=4\Pi_k+1$, which is the notes' map $a\mapsto4a+1$. So $\Pi_k\equiv5\pmod8$ for $k\ge1$. Every prime $P\equiv5\pmod8$ is solved by the shell $R=3$ with $u=2$. Hence only the first member of a predecessor chain can be a hard prime.

This is the classical Syracuse predecessor structure \cite{Lagarias}. It links this workbench with the Collatz workbench of \cite{WorkbenchReader}. No solvability statement iterates the Collatz map.
\end{remark}

\subsection{Single shells}\label{ssec:singleshells}

\begin{proposition}[non-coprime shells]\label{prop:noncoprime}
Let $p$ be prime and $R\equiv-p\pmod4$. Then $\gcd(R,pa)>1$ exactly when $p\mid R$; write $R=kp$.
\begin{enumerate}[label=(\alph*)]
\item If $p\nmid k$, the one-congruence criterion ``some $d\mid N^2$ has $d\equiv-N\pmod R$'' is equivalent to occupancy, and every such $d$ gives integral $y$ and $z$.
\item If $p^2\mid R$, it is not. At $p=73$, $R=7\cdot73^2$, the divisor $d=4\cdot73^3$ satisfies the congruence and gives $y=60$ but $z=1920/73$, and the shell has no solution.
\end{enumerate}
\end{proposition}
\status{Proved here; found by a reading pass. Four of the notes state the criterion without the coprimality hypothesis: the defect-completion article (Lemma~2.1), \emph{Split-zero residual moonshine} (Theorem~3.1), the $A_8$ snowflake note (Theorem~2.1) and \emph{Terminal-core lock covers} (Theorem~2.1). The error is harmless where the notes use it: the criterion can fail only when $p^2\mid R$, that is $R\ge3p^2$.}
\begin{proof}
(a) Put $m=(k+1)/4$, so that $a=pm$ and $N=p^2m$. A divisor $d$ with $d\equiv-N\pmod{kp}$ is divisible by $p$; write $d=pd_1$, so that $k\mid d_1+pm$.
\begin{itemize}
\item If $p^3\nmid d_1$, then $d_1\mid(pm)^2$ is a certificate for the reduced shell $(k,pm)$, and the formulas agree.
\item If $d_1=p^3w$ with $w\mid m^2$, then $k\mid4p^2w+1$, so $(w/k)=(-1/k)=-1$.
\item But every prime $\ell\mid m$ has $(\ell/k)=1$: for odd $\ell$ because $k\equiv-1\pmod\ell$ and reciprocity applies; for $\ell=2$ because $k\equiv7\pmod8$. So $(w/k)=1$, a contradiction.
\end{itemize}
(b) Direct computation, including a search over all $y$.
\end{proof}

\begin{proposition}[ramified shells]\label{prop:ramified}
Let $(R,p)$ be a coprime shell.
\begin{enumerate}[label=(\alph*)]
\item If $R\mid p-1$, then $E_a=M_a$.
\item If $R\mid p^2-1$, then $u\mapsto a^2/u$ maps $E_a$ to itself, and $|E_a|$ is odd exactly when $R\mid p+1$.
\end{enumerate}
\end{proposition}
\status{Proved here; checked on all $947$ coprime shells with $R\mid p^2-1$ and $p\equiv1\pmod4$ below $2000$.}
\begin{proof}
(a) $4a\equiv p\equiv1\pmod R$, so $4(u+a)\equiv4u+1$.

(b) The complement sends the class $-4^{-1}$ to $-4a^2\equiv-p^2/4$, which is $-4^{-1}$ again when $p^2\equiv1\pmod R$. The only possible fixed point is $u=a$. Finally $a\in E_a$ exactly when $R\mid4a+1=p+R+1$.
\end{proof}

\begin{proposition}[the pair-count series]\label{prop:pairzeta}
Let $A_R(n)$ be the number of triples $(X,Y,Z)$ with $XYZ=n$, $\gcd(X,Y)=1$ and $R\mid X+Y$. For a coprime shell, $A_R(pa)$ is the number of ordered certificates. For $R\ge3$,
\[
\sum_{(n,R)=1}\frac{A_R(n)}{n^s}=\frac1{\varphi(R)}\sum_{\chi\bmod R}\chi(-1)\,\frac{\zeta_R(s)\,L(s,\chi)\,L(s,\bar\chi)}{\zeta_R(2s)},
\]
where $\zeta_R$ omits the primes dividing $R$.
\end{proposition}
\status{Proved here. This is Theorem~4.3 of the defect-completion article; $8757$ coefficients were checked by a referee pass. No positivity statement follows.}
\begin{proof}
Write $1_{X\equiv-Y}=\varphi(R)^{-1}\sum_\chi\chi(-1)\chi(X)\bar\chi(Y)$. The variable $Z$ contributes $\zeta_R(s)$. Möbius inversion over the common divisors of $X$ and $Y$ contributes $L(s,\chi)L(s,\bar\chi)/L(2s,\chi\bar\chi)$, and $\chi\bar\chi$ is principal.
\end{proof}

The divisor-pair note asks for finitely many central and edge templates covering $S_{840}$. As a congruence cover this is impossible, because no template class contains a square modulo its modulus.
\begin{itemize}
\item For the edge templates this is Proposition~\ref{prop:nosquare}.
\item A central template is the class $n\equiv-R\pmod{4c}$ with $s\mid c^2$, $\gcd(c,R)=1$, $R\equiv3\pmod4$ and $R\mid s+c$. If $-R$ were a unit square modulo $4c$, then:
\begin{itemize}
\item every odd prime $\ell\mid c$ would have $(-R/\ell)=1$, hence $(\ell/R)=1$;
\item if $2\mid c$, then $-R\equiv1\pmod8$, so $(2/R)=1$.
\end{itemize}
So $(c/R)=(s/R)=1$. But $R\mid s+c$ gives $(s/R)=(-c/R)=-1$. (A reading pass found this argument.)
\end{itemize}
By Dirichlet's theorem, primes $p\equiv1$ modulo the product of $840$ and the template moduli exist, and they escape every template.

\subsection{Identities, locks and types}\label{ssec:locks}

\begin{proposition}[Rosati's family]\label{prop:rosati}
Let $k\ge1$, $c\mid k$ and $m=4k-1$. If $n\equiv-c\pmod m$, then
\[
\frac4n=\frac1{kn}+\frac1{k(n+c)/m}+\frac1{kn(n+c)/(mc)}.
\]
\end{proposition}
\status{Proved here; Salez's chart (14a) with $(B,C,D)=(c,1,k/c)$. The case $k=22$, $c=11$ solves every $n\equiv76\pmod{87}$. The two ``direct $29$-channels'' of \emph{Terminal residual coordinates} (Theorems~7.1--7.2), $n\equiv917\pmod{1276}$ and $n\equiv1605\pmod{2204}$, are correct. They are Salez's chart (15b) with $(B,C,F)=(1,29,11)$ and $(1,29,19)$. The first lies inside this family with $(k,c)=(80,40)$, $m=319$. The second lies inside chart (14a) with $(B,C,D)=(2,23,3)$, $m=551$. The notes' claim that these are the only direct channels holds only for lock divisors that are single primes from $\{5,11,19,23\}$. The lock with $\mu=5$ and $d=31$ solves $n\equiv429\pmod{620}$, a class containing the core prime $33{,}289$.}
\begin{proof}
Over the common denominator $kn(n+c)$ the three terms have numerators $n+c$, $mn$ and $mc$, which add up to $(1+m)(n+c)=4k(n+c)$. The denominators are integers because $m\mid n+c$ and $c\mid k$.
\end{proof}

\begin{proposition}[the full-divisor locks]\label{prop:locks}
Let $p$ be an odd prime, $\mu$ odd and $N=(p+\mu)/2$.
\begin{enumerate}[label=(\alph*)]
\item If $d\mid N$ is odd with $d\equiv-p\pmod{4\mu}$, then the shell $R=d$ carries the $E$-certificate $u=a/\mu$.
\item If $h\mid N$ with $h\equiv-1\pmod{4\mu}$, put $e=N/h$ and $t=(h+1)/(4\mu)$. Then $a=2\mu et$ and $R=\mu+2e$ satisfy $4a-p=R$, and $u=\mu^2t$ is an $M$-certificate.
\end{enumerate}
\end{proposition}
\status{Proved here. These are the locks of the notes, which a reading pass identifies with the faces $A=1$ of Type~I and $C=1$ of Type~II of \cite{ElsholtzTao}; lock (b) is Lopez's Type~B. In the notes' proof of (a) the numerator should read $ps+p+\mu$, not $ps+p+\mu s$.}
\begin{proof}
(a) Here $\mu\mid a$, and $4u+1=(p+d+\mu)/\mu$ is divisible by $d$.

(b) $4a-p=2e(h+1)-2eh+\mu=R$ and $u+a=\mu tR$.
\end{proof}

\paragraph{Computations on these certificates.}
\begin{itemize}
\item Every hard prime below $10^7$ has a lock certificate. A reading pass found the same below $10^8$, where the hard primes number $179{,}468$. Not every prime has one: $409\equiv1\pmod{24}$, which is not hard, has none for any odd $\mu$, by an exhaustive search. It is the only prime $p\equiv1\pmod4$ below $3\cdot10^6$ without one (found by the referee of this version).
\item Every prime $p\equiv1\pmod{24}$ below $10^7$ has a Type~II solution, that is, an $M$-certificate, in a shell $R\le107$. The record $p_*$ is the only one that needs $107$. A reading pass found the same below $10^8$.
\item The \emph{complement-square} certificates $u=Q^2$ with $Q\mid a$ (Lopez's Type~C) are not universal. The least prime $p\equiv1\pmod{24}$ without one is $97$, and the least hard prime without one is $3361$.
\end{itemize}

\subsection{Errors, and what the notes connect to}\label{ssec:noteserrors}

\paragraph{Errors.}
Besides the scope of the terminal core (Remark~\ref{rem:termcore}) and the unconditioned criterion (Proposition~\ref{prop:noncoprime}), the checked errors are the following.
\begin{itemize}
\item \emph{Constructive algebra of the centered Erdős--Straus resolution} states the conjecture (Theorems~15.1--15.3).
\begin{itemize}
\item Its Lemma~13.4 counts failure states as removed.
\item It uses an involution between two pieces of sizes $30{,}240$ and $37{,}800$, which cannot exist.
\item Its Theorem~15.2 treats the core as the residual locus.
\end{itemize}
The author's June replacement removed a conclusion of the same kind from its predecessor.
\item The converse of Star--Kneser forcing, stated in \emph{Star--Kneser positivity}, is false at $p_*$, $R=107$. It is already retracted in the archive.
\item The umbral group, the automorphism group of the glue code, is identified with $GL_2(3)$ in four notes. The argument does not exclude the binary octahedral group. The conclusion is true: the group has $13$ involutions (a reading pass).
\item The ``Coxeter--Fourier decomposition'' of the shell count over the six hard classes is ill-posed, because the count is not a function of $p\bmod840$. Already $A_{11}(p)$ takes the values $0,4,6,8,10$ on primes $p\equiv1\pmod{840}$.
\item The ``torsor map'' $(\ZZ/18)^\times\to S_{840}$ of the terminal papers is not a homomorphism.
\item The Kneser theorem used is that of \cite{Kneser}, of 1955, not the density paper of 1953.
\end{itemize}

\paragraph{Connections.}
\begin{itemize}
\item The ``explicit mixed mock functions'' of \emph{An explicit mixed-mock Virasoro support realization} (modulus $42$, labels $1,5,11$) are exactly Ramanujan's order-$7$ mock theta functions, in Zagier's vector $M_7$ \cite{Zagier}. They agree through $q^{40}$, up to one sign. For modulus $30$ one gets the order-$5$ functions. The first of them is also the invariant $q^{1/2}\widehat Z_0$ of the Brieskorn sphere $-\Sigma(2,3,7)$ \cite[eq.~(3.20)]{CFS}, in the framework of \cite{CCFGH}.
\item The note's identification of their shadow with the characters of the minimal model $M(2,7)$ fails. Its modified modular pair is not a representation of $SL_2(\ZZ)$, and no relabelling matches both $S$ and $T$.
\item The Busy-Beaver residues modulo $107$ in the Busy-Beaver notes are at chance level against a null model fixed in advance. The only logical link between the Busy-Beaver function and this equation is the classical one: a search for a counterexample is a $\Pi_1$ statement.
\item Every other link is correct where it was checked, with two exceptions, and none has a bearing on solvability. This covers Niemeier and umbral data, the class $6B$, Ogg's primes, Lorentzian and Apollonian geometry, and Cayley--Dickson algebras. The exceptions are:
\begin{itemize}
\item in the Lorentzian Descartes note, the ``modular flow'' is a rotation, not a Lorentz boost;
\item one isomorphism of the split-zero dossier (its Theorem~7.2) needs the relation $[\tau]\sim0$.
\end{itemize}
Both were found by reading passes.
\item The mechanisms that do bear on solvability are the classical ones: Elsholtz--Tao types \cite{ElsholtzTao}, Salez's charts \cite{Salez} and Lopez's types \cite{Lopez}.
\end{itemize}
